\item Pourquoi, Dieu, nous rejet\frflex{e}r sans fin ? Pourquoi cette colère sur les breb\frflex{i}s de ton troupeau ?
\item Rappelle-t\frflex{o}i la communauté que tu acqu\frflex{i}s dès l’origine,
\item La tribu que tu revendiqu\frflex{a}s pour héritage, la montagne de Sion où tu f\frflex{i}s ta demeure.
\item Dirige tes pas vers ces ru\frflex{i}nes sans fin, l’ennemi dans le sanctuaire a to\frflex{u}t saccagé ;
\item Dans le lieu de tes assemblées, l’advers\frflex{a}ire a rugi et là, il a plant\frflex{é} ses insignes.
\item On les a vus brand\frflex{i}r la cognée, comme en pl\frflex{e}ine forêt,
\item Quand ils bris\frflex{a}ient les portails à coups de m\frflex{a}sse et de hache.
\item Ils ont livré au fe\frflex{u} ton sanctuaire, profané et rasé la deme\frflex{u}re de ton Nom.
\item Ils ont dit : « All\frflex{o}ns ! Détruisons tout ! » Ils ont brûlé dans le pays les lie\frflex{u}x d’assemblées saintes.
\item Nos signes, nul ne les voit ; il n’y a pl\frflex{u}s de prophètes ! Et pour combien de temps ? Nul d’entre no\frflex{u}s ne le sait !
\item Dieu, combien de temps blasphémer\frflex{a} l’adversaire ? L’ennemi en finira-t-il de mépris\frflex{e}r ton Nom ?
\item Pourquoi reten\frflex{i}r ta main, cacher la f\frflex{o}rce de ton bras ?
\item Pourtant, Dieu, mon r\frflex{o}i dès l’origine, vainqueur des combats sur la f\frflex{a}ce de la terre.
\item C’est toi qui fendis la m\frflex{e}r par ta puissance, qui fracassas les têtes des drag\frflex{o}ns sur les eaux ;
\item Toi qui écrasas la t\frflex{ê}te de Léviathan pour nourrir les m\frflex{o}nstres marins ;
\item Toi qui ouvris les torr\frflex{e}nts et les sources, toi qui mis à sec des fle\frflex{u}ves intarissables.
\item À toi le jour, à t\frflex{o}i la nuit, toi qui ajustas le sol\frflex{e}il et les astres !
\item C’est toi qui fixas les b\frflex{o}rds de la terre ; l’hiver et l’été, c’est t\frflex{o}i qui les formas.
\item Rappelle-toi : l’ennemi a mépris\frflex{é} ton Nom, un peuple de fous a blasphém\frflex{é} le Seigneur.
\item Ne laisse pas la Bête égorg\frflex{e}r ta Tourterelle, n’oublie pas sans fin la v\frflex{i}e de tes pauvres.
\item Regarde vers l’Alliance : la gu\frflex{e}rre est partout ; on se cache dans les cav\frflex{e}rnes du pays.
\item Que l’opprimé éch\frflex{a}ppe à la honte, que le pauvre et le malheureux ch\frflex{a}ntent ton Nom !
\item Lève-toi, Die\frflex{u}, défends ta cause ! Rappelle-toi ces fous qui blasph\frflex{è}ment tout le jour.
\item N’oublie pas le vacarme que f\frflex{o}nt tes ennemis, la clameur de l’ennemi, qui m\frflex{o}nte sans fin.